% Pandoc template for the fed_fraud documents.
% Based on proposal_gen/templates/detailed_proposal_template/detailed_1_classic.tex
% (Classic report): centered title page, table of contents, ruled section
% headings and the Geodesic house footer. Fonts are Latin Modern (lmodern).
% Two layouts: the full report (title page and contents) and, with
% -V brief, a compact one-page-brief header in the same style.
\documentclass[letterpaper,11pt]{article}
\usepackage[T1]{fontenc}
\usepackage[utf8]{inputenc}
\usepackage{lmodern}
\usepackage[letterpaper,top=0.9in,bottom=1.2in,left=1in,right=1in,footskip=0.6in]{geometry}
\usepackage[table]{xcolor}
\usepackage{geodesicbrand}
\usepackage{amsmath,amssymb}
\usepackage{enumitem}
\usepackage{array}
\usepackage{calc}
\usepackage{longtable}
\usepackage{makecell}
\usepackage[most]{tcolorbox}
\usepackage{titlesec}
\usepackage{needspace}
\usepackage{etoolbox}
\usepackage{hyperref}
\hypersetup{colorlinks=true,urlcolor=coolblack,linkcolor=coolblack,filecolor=geoblue,
  pdftitle={Treasury (IRS and Fiscal Service)},pdfauthor={\geofirm}}
\urlstyle{same}

\setlist{itemsep=2pt,leftmargin=1.4em}
\setlength{\parindent}{0pt}
\setlength{\parskip}{6pt}
\providecommand{\tightlist}{\setlength{\itemsep}{0pt}\setlength{\parskip}{0pt}}
\setcounter{secnumdepth}{0}

% Section style: house rule under each heading
\titleformat{\section}{\Large\bfseries\color{coolblack}}{\thesection}{0.6em}{}[{\color{coolblack}\titlerule[0.5pt]}]
\titleformat{\subsection}{\large\bfseries\color{geomid}}{\thesubsection}{0.6em}{}
\titleformat{\subsubsection}{\normalsize\bfseries\color{geomid}}{\thesubsubsection}{0.6em}{}
\titlespacing*{\section}{0pt}{10pt}{6pt}

% Tables are written by latex/filters.lua (navy header, banded rows)
\setlength{\LTpre}{6pt}\setlength{\LTpost}{6pt}

\begin{document}

% BRIEF HEADER
\thispagestyle{geodesic}
\begin{minipage}[c]{0.13\textwidth}\includegraphics[width=\linewidth]{icon.png}\end{minipage}\hfill
\begin{minipage}[c]{0.84\textwidth}
  {\textsc{\large \geofirm}}\par\vspace{2pt}
  {\small\color{geogray}Federal fraud expansion\enspace\textperiodcentered\enspace When: \compiletime}
\end{minipage}\par\vspace{6pt}
{\color{coolblack}\rule{\linewidth}{0.3mm}}\par\vspace{2pt}
{\LARGE\bfseries\color{coolblack} Treasury (IRS and Fiscal Service)}\par
\vspace{2pt}{\large\color{geomid} One-page brief}\par
\vspace{-2pt}{\color{coolblack}\rule{\linewidth}{0.3mm}}\par\vspace{4pt}
\pagestyle{geodesic}

\textbf{Priority:} 8 of 33 (score 3.5). \textbf{FY2026 outlays:}
\$1,975B (USAspending; mostly interest on the debt). \textbf{As of:}
2026-10-27. Public data only.

\hypertarget{bottom-line}{%
\section{Bottom line}\label{bottom-line}}

Treasury holds the largest single fraud-prone payment outside health
care: refundable tax credits paid through the IRS. The Earned Income Tax
Credit alone had \$21.1B in improper payments in FY2025, and the
pandemic Employee Retention Credit drew billions in fraudulent claims.
Taxpayer data are confidential, so public work has to start from
geography, preparers and promoters rather than individual returns. The
payoff is large, but this is a proposal-first agency.

\hypertarget{exposure}{%
\section{Exposure}\label{exposure}}

\begin{itemize}
\tightlist
\item
  EITC: \$21.1B improper (32.7\%). Additional Child Tax Credit: 13.6\%
  rate. American Opportunity Tax Credit: 30.7\%. Premium Tax Credit:
  31.6\%. (GAO, FY2025)
\item
  Employee Retention Credit: IRS Criminal Investigation had opened 545
  ERC cases involving more than \$5.6B by February 2025, out of 2,039
  COVID fraud investigations totaling \$10B.
\item
  Fiscal Service disburses most federal payments and runs Do Not Pay,
  which makes it the natural home for any cross-agency screen.
\end{itemize}

\hypertarget{public-data-we-can-use}{%
\section{Public data we can use}\label{public-data-we-can-use}}

\begin{itemize}
\tightlist
\item
  \textbf{IRS Statistics of Income (SOI) ZIP code data:} EITC and
  refundable credit counts and dollars by ZIP, by year.
\item
  \textbf{IRS Return Preparer data and enforcement releases:} preparer
  injunctions and DOJ Tax Division cases name preparers and promoters.
\item
  \textbf{PPP loan-level data (SBA):} employer names, addresses and jobs
  reported, which can be set against public ERC promoter cases.
\item
  \textbf{IRS-CI and DOJ press releases} as labeled examples.
\end{itemize}

\hypertarget{detection-plays}{%
\section{Detection plays}\label{detection-plays}}

\begin{enumerate}
\def\labelenumi{\arabic{enumi}.}
\tightlist
\item
  \textbf{Credit hot spots:} ZIP codes where EITC or ACTC claim rates
  jump far above what income and household data predict, flagged for
  preparer review.
\item
  \textbf{Promoter networks:} link charged ERC and EITC promoters to the
  businesses, addresses and preparers that recur across cases.
\item
  \textbf{Pandemic overlap:} businesses with PPP loans that also appear
  in ERC cases, to estimate double dipping and find similar unreviewed
  firms.
\end{enumerate}

\hypertarget{first-prototype}{%
\section{First prototype}\label{first-prototype}}

An EITC hot-spot map from SOI ZIP data, with a ranked list of anomalous
areas and the known preparer cases inside them. Validate against DOJ
preparer injunctions.

\hypertarget{limits}{%
\section{Limits}\label{limits}}

Returns, preparer identifiers on returns, and ERC claims are
confidential (26 U.S.C. 6103). Public data can only point to places and
promoters; return level scoring needs an IRS partnership.

\hypertarget{who-to-take-it-to}{%
\section{Who to take it to}\label{who-to-take-it-to}}

IRS Return Preparer Office and Criminal Investigation, TIGTA, Treasury's
Office of Payment Integrity (Do Not Pay), and DOJ Tax Division.

\hypertarget{sources}{%
\section{Sources}\label{sources}}

\begin{itemize}
\tightlist
\item
  \href{https://www.gao.gov/products/gao-26-108694}{GAO-26-108694:
  FY2025 improper payments}
\item
  \href{https://irs.gov/compliance/criminal-investigation/five-years-post-cares-act-irs-ci-has-launched-2039-covid-fraud-investigations-totaling-10b-in-attempted-fraud}{IRS-CI:
  2,039 COVID fraud investigations, 545 ERC cases (March 2025)}
\item
  \href{https://www.crfb.org/blogs/federal-improper-payments-total-186-billion-fy-2025}{CRFB:
  FY2025 improper payments by program}
\item
  \href{https://www.usaspending.gov/agency}{USAspending agency profiles
  (FY2026 outlays)}
\end{itemize}

\end{document}
